\documentclass[border = 1mm]{standalone}
% Caption:
% Compton Scattering: gamma + e -> gamma + e.
% Reproduction from Fig. 6.12 of Halzen & Martin "Quarks
% and Leptons: An Introductory Course in Modern Particle
% Physics".
%
% Autor:
% R. Sánchez-Martínez, Jun 2026.
%
% Figure published at:
% instagram.com/ars_mathematica
%
% Generated with TikZ and compiled preferable with XeLaTeX.

%
\usepackage[utf8]{inputenc}
\usepackage[T1]{fontenc}
\usepackage{microtype}
%
\usepackage{amsmath,amssymb,amsfonts}
\usepackage{slashed}
%
\usepackage{xcolor}
\usepackage{tikz}
\usetikzlibrary{
	babel,
	math,calc
}
\usepackage{tikz-feynman}

\usepackage{fontspec}
\usepackage[default]{fontsetup}

\begin{document}
% -- Figure
% Key concepts:
%	— \draw[style] (a) to (b);
%	  draws a line from (a) to (b) with a specific
%     graphic [style].
%   — \draw (a) to node {text} (b);
%	  adds a node containing {text} in the midpoint
%     of (a) and (b).
%   — ... node[anchor = base] {text} ...
%     indicates that the coordinate specification will
%     refer to the node's baseline midpoint, useful
%     for aligning nodes of different heights.
\begin{tikzpicture}[
	line join = round,
	line cap  = round,
	line width = 0.62pt,
]

	% Main box 
	\node[
		inner sep = 0pt,
		minimum width  = 6cm,
		minimum height = 8cm,
		name = N,
	] at (0,0) {};
	
	% Clipping mask, it deletes everything outside the box (N),
	% useful for obtaining well-defined-size drawings
	\clip (N.north west) rectangle (N.south east);

	% Geometric definitions
	\tikzmath{
		\d = 2.5;
		\l = 1.4; \th = 28;
	}

	% Labels and equations
	\node[align = center] (title)
		at (0,3)
		{Dispersión Compton};
		
	\node[inner sep = 1mm, draw = gray!75, yshift = -3mm]
		at (title.south)
		{$\gamma + e^- \to \gamma + e^-$};
		
	\node[anchor = base, font = \small]
		at (0,-2.7)
		{$
			-e^2\bar{u}(p'){\epsilon'_\nu}^{*}\gamma^\nu
			\dfrac{\slashed{p}+\slashed{k}+m}{(p + k)^2 - m^2}
			\gamma^\mu\epsilon_\mu u(p)
		$};

	% ---
	\begin{scope}[
	% use shift to center the picture or placing it
	% wherever you wan
		shift = {(-0.5*\d,-0.2)},
	]
	
	% Main drawing using TikZ-Feynman macros
	\begin{feynman}
	
		% coordinates of the vertices
		\coordinate (v1) at (0,0);
		\coordinate (v2) at (\d,0);
		
		\coordinate[shift = {(-90-\th:\l)}] (p) at (v1);
		\coordinate[shift = {(+90+\th:\l)}] (k) at (v1);
		
		\coordinate[shift = {(-90+\th:\l)}] (pp) at (v2);
		\coordinate[shift = {(+90-\th:\l)}] (kk) at (v2);
		
		% main lines drawings
		\draw[fermion]
			(v1)
			to
			node[anchor = base, shift = {(0,0.6)}, font = \small]
				{$\displaystyle
					\frac{i\sum u\bar{u}}{(p + k)^2 - m^2}
				$}
			(v2);
		
		\draw[fermion]
			(p) node[anchor = base, shift = {(-0.05,-0.3)}] {$p$} 
			to node[anchor = base, shift = {(-0.3,0)}] {$u$} (v1);
		\draw[fermion]
			(v2) to node[anchor = base, shift = {(+0.3,0)}] {$\bar{u}$} (pp)
			node[anchor = base, shift = {(+0.10,-0.3)}] {$p'$};
		\draw[photon]
			(k)
			node[anchor = base, shift = {(0,+0.15)}] {$k$} 
			to
			node[anchor = base, shift = {(-0.35,0)}] {$\epsilon_\mu$}
			(v1);
		\draw[photon]
			(v2)
			to
			node[anchor = base, shift = {(+0.45,0)}] {${\epsilon'_\nu}^*$}
			(kk)
			node[anchor = base, shift = {(+0.15,+0.15)}] {$k'$};
		
		% dots
		\fill (v1) circle (1.2pt);
		\fill (v2) circle (1.2pt);
		
	\end{feynman}
	\end{scope}
	% ---
	
\end{tikzpicture}
\end{document}